\section{Conclusion}
\label{sec:conclusion}

This work systematically evaluates Generative AI's dual impact on spam detection, revealing fundamental asymmetries in offensive and defensive capabilities. Our findings demonstrate three critical insights: (1) LLM-generated synthetic spam can partially substitute real training data, though effectiveness varies dramatically by model generation quality; (2) AI-powered spam poses moderate detection challenges with performance strongly influenced by classifier architecture and training strategies; (3) cross-model synthetic augmentation significantly improves detection, with diversity-driven approaches outperforming quality-focused methods. Ensemble classifiers consistently demonstrate superior robustness to distributional shifts and cross-model generalization compared to linear models.

For cybersecurity practitioners, these results indicate that Generative AI currently favors attackers through accessible generation capabilities and low deployment costs. However, defenders can maintain effectiveness by leveraging robust classifier architectures, cross-model training augmentation, and strategic synthetic data integration. Our statistical frameworks enable practitioners to predict detection performance across arbitrary synthetic mixing ratios, supporting evidence-based deployment decisions.

Future research should explore open-source LLMs to assess democratization effects, investigate transformer-based detection architectures, and evaluate generalization to broader cybersecurity domains including phishing, social engineering, and adversarial content generation.
